% !TEX root = ../main/main.tex
% chapter_1_intro.tex  
\chapter{Introduction and Motivation}
\label{chap:intro}

\section{Why Focus on the NFL}

The \textbf{National Football League (NFL)} is an ideal testbed for quantitative decision systems because it combines: (i) deep liquid betting markets; (ii) abundant, granular data; and (iii) inherently sequential strategic decisions. These properties support rigorous modeling and credible evaluation of edge.

\begingroup\RaggedRight\sloppy\hbadness=10000
\begin{itemize}
  \item \textbf{Liquidity.} Many books and large volumes enable precise price comparisons and operationally meaningful CLV measurement.
  \item \textbf{Data richness.} Public play-by-play, injury reports, and weather feeds support feature engineering and validation at weekly cadence.
  \item \textbf{Sequential decisions.} Weekly cycles and interdependent markets (spread/total/SGP) invite RL-style policies with risk controls.
\end{itemize}
\endgroup

\section{Research Questions and Objectives}

We pursue four questions:
\begin{enumerate}
  \item \textbf{Architecture.} How to combine classical models, modern ML, and RL coherently for NFL markets?
  \item \textbf{Uncertainty \& risk.} How to quantify uncertainty and translate it into safe stake sizing?
  \item \textbf{Value of information.} What is the marginal lift from feature families (injury, rest, weather, microstructure)?
  \item \textbf{Evaluation.} How to simulate and measure policies credibly under realistic frictions and dependence?
\end{enumerate}

\noindent Deliverables include a unified modeling stack, a risk-aware staking module, a simulator for strategy evaluation, and a reproducible system-of-systems.

% (Contributions merged into Research Questions; no separate section)

\section{Scope and Boundaries}
We target pre-game markets across two domains: (1)~\textbf{game-level} (spread/total/moneyline) from 1999 onward, and (2)~\textbf{player props} (passing yards, rushing yards, receptions) from 2020 onward. Live in-game betting is out of scope. Only public data are used. All evaluation is out-of-sample with rolling-origin splits and clear decision-time information. Critically, we find that \textbf{performance varies dramatically by market type}: spreads show $-7.5$\% ROI (market efficient) while props showed $+2.96$\% ROI in a limited test window (pipeline currently non-functional).

\section{Thesis Statement and Hypotheses}

This dissertation investigates a fundamental question in quantitative sports analytics: \textit{Can rigorous statistical methods and machine learning extract systematic profits from NFL betting markets using publicly available data?}

We develop a complete betting system pipeline—from data ingestion to model training to risk management to evaluation—and test it rigorously on 5,529 games across 21 seasons (2004--2024).

\textbf{Thesis.} A disciplined ensemble architecture combining orthogonalized features, diversified model methodologies, and rigorous validation reveals \textbf{differential market efficiency across bet types}. NFL spread betting markets are efficient for publicly available information, but \textbf{player props markets show exploitable edges}. Most surprisingly, model complexity does not guarantee performance improvement.

\textbf{Critical findings (December 2025):} Performance varies dramatically by market type:
\begin{itemize}[itemsep=2pt,parsep=0pt]
    \item \textbf{Spread betting:} $-7.5$\% ROI, 51.0\% win rate---market efficient, vigorish barrier cannot be overcome
    \item \textbf{Player props:} $+2.96$\% ROI, 53.89\% win rate in limited test window---pipeline currently non-functional (model files missing, data tables empty)
    \item \textbf{Simple power ratings:} 57.7\% ATS, +10.2\% ROI on 2025 ($N\!=\!221$); multi-year 2020--2025: 51.8\% ATS, $-1.1$\% ROI (1,310 games)---single-season results misleading
    \item \textbf{Position stability matters:} QB passing (0.85 YoY correlation) vs RB rushing (0.50)---not all props are equal
\end{itemize}
The conclusion is nuanced: spread markets are efficient, single-season results can be misleading (power ratings multi-year ATS is 51.8\%), and props pipeline requires restoration before claims can be validated. Prediction accuracy does not imply betting profitability, and honest multi-year evaluation matters more than model sophistication.

\textbf{Supporting Evidence:}
\begin{enumerate}[itemsep=2pt,parsep=0pt]
    \item \textbf{Orthogonalization principle} (Appendix~\ref{app:orthogonalization}): Two-stage feature orthogonalization architecture separates player skill from opponent strength via Gram-Schmidt decomposition. Initial results showed promise (88\% in-sample improvement) but rigorous out-of-sample validation revealed target leakage---a critical lesson documented in Appendix~N demonstrating that validation rigor is paramount in prediction systems.

    \item \textbf{Ensemble diversification} (Appendix~\ref{app:architecture-evolution}): 5-way architecture (State-Space 27\%, Hierarchical 27\%, XGBoost 23\%, Informative Priors 13\%, Hybrid 10\%) with mean correlation $\bar{\rho} = 0.32$ achieved 56\% variance reduction. Ensemble variance decomposition principle: Adding uncorrelated weak models beats refining correlated strong models.

    \item \textbf{Hybrid architectures} (Chapter~\ref{chap:lessons-learned}): 70\% static + 30\% incremental XGBoost blend achieved +4.2\% improvement over pure static (38.03 vs 39.66 MAE average 2022--2024), resolving static-adaptive dilemma through complementary strength balancing.

    \item \textbf{Multi-season validation} (Appendix~\ref{app:orthogonalization}, Table~\ref{tab:multiseas-validation}): Three-season consistency testing (2022--2024, 546 games) with inter-season variance $\sigma = 0.40$ MAE (vs baseline $\sigma = 1.89$ MAE) prevented overfitting to single-season quirks.

    \item \textbf{Honest failure documentation} (Chapter~\ref{chap:results}, \S\ref{sec:failure-analysis}): Transparent accounting of Bayesian Neural Network investigation (84 hours, 6 failed approaches, 0 production models) and baseline error (incremental vs mean predictor, 9 hours wasted) with Bayesian abandonment criteria ($k = 2$ failures $\to$ terminate) prevents repeat mistakes.
\end{enumerate}

\subsection{Methodological Contributions}

This work makes six primary contributions, independent of profitability outcomes:
\begin{enumerate}
  \item \textbf{Feature orthogonalization methodology ("Ridge Residuals"):} Formalized the opponent-player confound as a Gram-Schmidt orthogonalization problem (Appendix~\ref{app:orthogonalization}), proving that residual extraction eliminates spurious correlations (Theorem~\ref{eq:orthogonal-decomposition-theorem}). The three-stage Ridge Residuals architecture (Ridge opponent baseline $\to$ QB-specific residual extraction $\to$ XGBoost with uncorrelated features) generalizes to any domain with confounding (NBA player-defender, MLB pitcher-park, soccer striker-defense, finance market-neutral alpha). Initial results showed promise (88\% in-sample improvement), but rigorous out-of-sample validation revealed target leakage in the feature pipeline---a critical lesson documented in Appendix~N demonstrating that validation rigor is paramount and in-sample improvements can be misleading. The methodology remains theoretically sound and applicable to future work with proper temporal splits.

  \item \textbf{Ensemble architecture evolution methodology}: Documented complete 13-month journey from v1.0 catastrophic failure ($-83.5\%$ drawdown) through multiple iterations to v6.0 production system, culminating in comprehensive data leakage investigation that confirmed \textbf{market efficiency} (cleaned baseline: 51.0\% win rate, $-7.5\%$ ROI). Contributions: (a) ensemble variance decomposition analysis proving diversification dominates individual sophistication ($\bar{\rho}$ reduction more cost-effective than M increase), (b) domain-specific weight optimization (Dirichlet for props, static for spreads), (c) hybrid static-incremental architecture pattern (70/30 blend resolves stability-adaptation tradeoff), (d) honest failure accounting framework with Bayesian abandonment criterion (prevents sunk cost fallacy), (e) \textbf{rigorous leakage detection} that caught Ridge Residuals target leakage (88\% ``improvement'' $\to$ 47\% degradation) and CQL target contamination (100\% win rate $\to$ DISABLED).

  \item \textbf{Production deployment playbook with multi-layer validation}: Documented 6-layer defense-in-depth architecture (Chapter~\ref{chap:results}, Table~\ref{tab:defense-in-depth-condensed}) catching different error classes at different stages. Key insight: Baseline error (incremental vs mean predictor) passed unit tests and integration tests but failed at Layer 3 (multi-season historical validation), preventing production disaster. Framework includes: fail-fast validation on critical dependencies, type-safe model loading for polyglot systems (R/Python integration), canonical training standards (2018--2024 + H=4 decay), and multi-season consistency gates (minimum 3 seasons, $\sigma < 2$ yards inter-season variance).

  \item \textbf{Rigorous empirical evaluation}: Comprehensive testing of weather effects, calibration-to-profitability gaps, and RL vs Kelly baselines with transparent reporting of negative results.
  \item \textbf{Complete system architecture}: A full pipeline from ingestion to evaluation, reusable for sports analytics and sequential decision domains.
  \item \textbf{Dependence-aware methods}: Copula-based pricing for correlated outcomes (same-game parlays, teasers) and key-number reweighting for integer margins.
  \item \textbf{Transparent governance}: Documentation of model selection, risk constraints, and operational decision logic, including when the system declines to act.
\end{enumerate}

\noindent The dissertation demonstrates that strong calibration, while necessary, is insufficient for profitability in semi-strong efficient markets where vigorish imposes a 4--5\% hurdle. Full results appear in \Cref{chap:results}.

\section{Reproducibility and Ethics}
Pipelines are containerized and deterministic. Seeds, dataset manifests, and artefacts are versioned. We adopt responsible-gambling principles with exposure caps, volatility limits, and stop-losses; results are presented with uncertainty, not as guarantees.

\section{Chapter Summary}
This chapter motivated the NFL as a fertile ground for rigorous decision analytics and laid out the claims and scope of the work. The remainder of the dissertation builds from data foundations (\Cref{chap:data}) through baselines (\Cref{chap:methods}), risk management (\Cref{chap:risk}), and simulation (\Cref{chap:sim}) to a deployable policy.

\section{Dissertation Structure}

Below is a brief road map of this dissertation:

\begin{itemize}
  \item \textbf{Chapter~2: Literature Review} — survey of classical and modern models in sports prediction, betting markets, RL in game domains.
  \item \textbf{Chapter~3: Data Foundations and Feature Engineering} — discussion of NFL data sources, structure, preprocessing, feature catalogs, era handling.
  \item \textbf{Chapter~4: Baseline Models} — implementation and calibration of GLM, state-space, Poisson, and classical benchmarks.
  \item \textbf{Chapter~5: Advanced Bayesian Methods} — hierarchical models, ensemble stacking, Dirichlet optimization, GPU-accelerated inference.
  \item \textbf{Chapter~6: Reinforcement Learning Framework} — state/action/reward specification, offline RL design (CQL, IQL), training pipelines.
  \item \textbf{Chapter~7: Uncertainty \& Risk Management} — posterior distributions, risk-aware betting, Kelly strategies, CVaR constraints, drawdown analysis.
  \item \textbf{Chapter~8: Simulation \& Strategy Testing} — Monte Carlo engines, teasing, parlays, correlated outcomes, strategy performance.
  \item \textbf{Chapter~9: Results, Ablations, Discussion} — comparative performance, feature ablations, robustness, error analysis.
  \item \textbf{Chapter~10: Production Deployment} — modular pipeline, experiment tracking, deployment strategy, monitoring, version control.
  \item \textbf{Chapter~11: Lessons Learned} — synthesis of 13-month research journey from catastrophic failure to market efficiency confirmation. Technical lessons, honest failure archive (BNN investigation, Ridge Residuals leakage, CQL target contamination), statistical insights, and transferable applications.
  \item \textbf{Chapter~12: Conclusion and Future Work} — final synthesis, limitations, and research frontiers.
  \item \textbf{Appendices} — extended figures, proofs, code reference, glossaries.
  \item \textbf{Appendix~N: Orthogonalization Deep Dive} — Mathematical treatment of Ridge Residuals methodology (Gram-Schmidt orthogonalization, variance decomposition, omitted variable bias), ablation studies, multi-season validation (2022--2024 consistency analysis), production deployment lessons (fail-fast, defense-in-depth), and future directions (nonlinear opponent baseline, hierarchical residuals, causal features).
  \item \textbf{Appendix~M: Architecture Evolution} — Complete timeline from v1.0 catastrophic failure ($-83.5\%$ drawdown, October 2024) to v6.0 production system (November 2025), culminating in December 2025 data leakage investigation confirming market efficiency (cleaned baseline: 51.0\% win rate, $-7.5\%$ ROI). Honest failure documentation (BNN abandonment, Ridge Residuals leakage, CQL target contamination), convergence improvement analysis (6000-iteration MCMC optimization), CVX weight optimization validation, hybrid architecture development (70/30 static-incremental blend), and lessons learned framework.
  \item \textbf{Appendix~O: Model Specifications} — Complete technical specifications for 6 production models (State-Space v2, Hierarchical v2, XGBoost v4 Ridge Residuals \textbf{[INVALIDATED---target leakage]}, Informative Priors v2, Hybrid XGBoost v1, Supporting models). Includes: MCMC convergence diagnostics ($\hat{R}$, ESS, divergences), hyperparameter configurations, multi-season validation tables (2022--2024 performance), feature importance analysis, computational benchmarks, software version registry (reproducibility standard), random seed registry, and 12-item deployment checklist.
\end{itemize}

\section{Technical Approach Overview}

At a high level, the system will operate in layers:

\begin{itemize}
  \item A data ingestion and feature pipeline — building situational, team-level, market-level features.
  \item Classical and benchmark models (GLM, Poisson, state-space) to form priors and baselines.
  \item An ensemble prediction system combining diversified methodologies (Bayesian state-space, hierarchical partial pooling, gradient boosting, ELO ratings), and adaptive weight optimization (CVX convex programming, Weighted Least Squares with exponential recency decay). \textbf{Note:} Ridge Residuals orthogonalization showed promise but was \textbf{invalidated} due to target leakage (see Appendix~N). Architecture evolution from v1.0 single-model catastrophe to v6.0 5-way ensemble detailed in Appendix~\ref{app:architecture-evolution}.
  \item Conservative offline RL agent (CQL with posterior-variance gating) for bet sizing, designed for off-policy evaluation with high-confidence bounds. \textbf{Critical note:} CQL is currently \textbf{DISABLED} pending retraining after target leakage was discovered in the training pipeline (see ADR-016). Initial claims of +24\% improvement were invalidated; simpler Kelly-LCB baselines are used until CQL is retrained on clean data.
  \item Uncertainty propagation (posterior distributions, bootstrap ensembles) to inform risk control.
  \item A simulation engine that translates distributions into actionable betting strategies.
  \item An evaluation and governance layer to compare models, version them, log experiments, and deploy.
\end{itemize}



\section{Glossary of Key Terms}

Here are some terms you’ll see frequently:
\begin{description}
  \item[CLV:] Closing Line Value — difference between market-implied and model-implied edge.
  \item[Kelly Fraction:] Optimal fraction of bankroll to stake given edge and odds.
  \item[Posterior Uncertainty:] Bayesian credible intervals on model predictions.
  \item[Middling / Teasers:] Betting strategies that exploit distributions across markets.
  \item[Feature ‹class›:] A group of inputs, e.g.\ injuries, rest, market signals.
  \item[Ensemble:] A weighted combination of predictions from multiple models.
\end{description}

% (Removed Writing and Style Notes)

\chaptersummary{
This chapter established the motivation, scope, and core contributions of the dissertation. We test whether NFL betting markets can be beaten with publicly available data. The comprehensive data leakage investigation (December 2025) reveals \textbf{market efficiency for spreads}: $-7.5$\% ROI. Player props showed $+2.96$\% ROI in limited testing but the pipeline is non-functional. A simple power ratings model showed 57.7\% ATS in 2025 but $-1.1$\% ROI over 6 years. The contributions are methodological (rigorous validation, leakage detection, honest reporting) and cautionary (single-season results can be misleading; multi-year evaluation is essential).
}{
We now survey the foundational literature—canonical models (Harville, Stern, Glickman), score distributions (Poisson, Skellam, bivariate), evaluation frameworks (proper scoring, calibration), and market microstructure—in \Cref{chap:litreview}.
}
